\begin{abstract}
Хоёр болон түүнээс дээш хүний хамтын хөдөлгөөний синхрончлол нь хос болон бүлгийн бүжиг, спортын хос болон багийн төрөл, дасгалжуулагч ба тамирчны дасгал, сэргээн засах эмчилгээ зэрэг салбарт гүйцэтгэлийн чанарыг тодорхойлох чухал үзүүлэлт юм. Уламжлалт байдлаар синхрончлолыг дасгалжуулагч, шүүгч эсвэл мэргэжилтэн нүдээр ажиглан үнэлдэг бөгөөд энэ үнэлгээ нь хүний субъектив дүгнэлтээс ихээхэн хамаардаг, тоон хэмжүүргүй, давтагдах боломж муутай байдаг. Мөн хөдөлгөөн бүртгэх (motion capture) тоног төхөөрөмж ашиглах нь өндөр өртөгтэй, тусгай мэдрэгч зүүх шаардлагатай тул өдөр тутмын хэрэглээнд тохиромжгүй. Иймд хүний байрлал тодорхойлох (Human Pose Estimation, HPE) технологи ашиглан хүмүүсийн биеийн гол цэгүүдийг тодорхойлж, тэдгээрийн хөдөлгөөний ижил төстэй байдал болон синхрончлолын түвшнийг тоон утгаар үнэлэх арга боловсруулах нь онолын болон практик ач холбогдолтой асуудал болж байна.

\medskip\noindent\textbf{Судалгааны ажлын зорилго}\par
Энэхүү судалгааны ажлын зорилго нь хоёр болон түүнээс дээш хүний хөдөлгөөний өгөгдөл өгөгдсөн үед HPE загвар ашиглан тэдгээрийн биеийн гол цэгүүдийг тодорхойлж, гол цэгүүдийн байрлал болон үений өнцөгт үндэслэн хүмүүсийн хөдөлгөөний ижил төстэй байдал болон синхрончлолын түвшнийг тоон хэмжүүрээр тооцоолох алгоритм боловсруулах явдал юм. Хэрэгжүүлэлтийг хийхдээ үндсэн тохиолдолд хоёр хүний хөдөлгөөнийг харьцуулах бөгөөд хүний тоо нэмэгдэхэд хүмүүсийг хос хосоор нь харьцуулж, дундажлан бүлгийн синхрончлолыг үнэлнэ. Мөн боловсруулсан алгоритмын үр дүнг синхрончлолын түвшин урьдчилан мэдэгдэж буй туршилтын өгөгдөл дээр шалгаж нарийвчлал болон гүйцэтгэлийг туршилтаар үнэлэх, практикт ашиглах боломж болон хязгаарлалтыг тодорхойлоход оршино.

\medskip\noindent\textbf{Судалгааны ажлын зорилтууд}\par
Тус зорилгыг хэрэгжүүлэхийн тулд дараах зорилтуудыг дэвшүүлэв. Үүнд:

\begin{enumerate}
\def\labelenumi{\arabic{enumi}.}
\item
  Хүний байрлал тодорхойлох технологийн онол, арга зүй болон хөдөлгөөний синхрончлолыг үнэлэх одоо ашиглагдаж буй аргачлалуудыг судлах
\item
  Видео бичлэгээс хүний биеийн гол цэгүүдийг илрүүлэх HPE загваруудыг (OpenPose, MediaPipe, YOLO-Pose, HRNet гэх мэт) харьцуулан судалж, шаардлагатай гол цэгийн тоо болон гүйцэтгэлд үндэслэн туршилтад ашиглах загварыг сонгох
\item
  Хоёр болон түүнээс дээш хүний хөдөлгөөнийг агуулсан туршилтын өгөгдлийг бэлтгэж, бичлэгт тавигдах шаардлагыг тодорхойлох
\item
  Илэрсэн гол цэгүүдийн өгөгдлийн дутуу цэгийг нөхөх, шуугиан арилгах урьдчилсан боловсруулалтын алгоритм боловсруулах
\item
  Гол цэгүүдийн байрлал, үений өнцөг болон хугацааны цуваанд үндэслэн хүмүүсийн хөдөлгөөний ижил төстэй байдал болон бүлгийн синхрончлолын түвшнийг тооцоолох алгоритм боловсруулах
\item
  Боловсруулсан алгоритмыг туршиж, синхрончлолын үнэлгээний үр дүнг синхрончлолын түвшин урьдчилан мэдэгдэж буй бичлэгүүд дээр шалгаж нарийвчлал, боловсруулалтын хурд болон найдвартай байдлыг статистик хэмжүүрээр үнэлэх
\item
  Видео бичлэгийн нягтаршил, кадрын давтамж зэрэг оролтын өгөгдлийн шинж чанар алгоритмын үр дүнд хэрхэн нөлөөлөхийг туршилтаар судалж, аргын боломж болон хязгаарлалтыг тодорхойлох
\end{enumerate}

\medskip\noindent\textbf{Судалгааны үр дүн}\par
Таван HPE загварыг харьцуулж RTMPose-m загварыг сонгон, үений өнцгийн зөрүү, корреляц, хөндлөн корреляцаар олсон цагийн хоцрогдлыг нэгтгэсэн синхрончлолын алгоритмыг вэб хөтөч болон сервер дээр ажилладаг систем болгон хэрэгжүүлэв. Нэг хүний бичлэгээс FFmpeg-ээр үүсгэсэн 32 хяналттай бичлэг дээр алгоритм 0--1 секундийн бүх хоцрогдлыг 0 мс алдаатай олж, синхрон (99.8), хоцрогдолтой (54.3--97.1) болон өөр хөдөлгөөнтэй (45.4--62.0) бичлэгүүдийг ялгаж, толин тусгал хөдөлгөөнийг таньж, 3 ба 4 хүнтэй бичлэгт хоцорсон хүнийг зөв илрүүлэв. Хүний хэмжээ 50\%, нягтаршил 360p, кадрын давтамж 10 кадр/с хүртэл буурахад үр дүн бараг өөрчлөгдөөгүй бол хүний өндөр ойролцоогоор 140 пикселээс бага үед хүний тоог автоматаар тоолох алхам алдаа гаргасан нь аргын гол хязгаарлалт болов.
\end{abstract}
